\section{Chương IV: Kết luận}
\\ Mô hình đã được xây dựng nhằm phân loại ung thư vú thành hai loại: ác tính (malignant) và lành tính (benign) dựa trên các đặc trưng trong bộ dữ liệu Breast Cancer Wisconsin. Bộ dữ liệu này bao gồm các thông tin về đặc tính tế bào học của khối u, như là kích thước hạt nhân, hình dạng và kết cấu.
\\ Bằng cách sử dụng Google Colab, tổng thời gian chạy của mỗi thuật toán là khoảng dưới 1 phút. Độ chính xác (accuracy) của các mô hình được sử dụng trong cuộc điều tra này đều ở mức khá cao trên 92\%, cho thấy những mô hình này có độ tin cậy nhất định. 
\\ Ngoài độ chính xác (accuracy), chỉ số recall cũng là nhân tố quan trọng trọng việc lựa chọn mô hình phù hợp. Từ kết quả thực thi đã có, ta thấy mô hình Logistic Regression có độ chính xác (accuracy) và chỉ số recall cao nhất nên mô hình này phù hợp nhất với yêu cầu bài toán và bộ dữ liệu cho trước.  
\\Một số đề xuất để cải thiện mô hình dự đoán:
\begin{itemize}
    \item Thử nghiệm thêm các phương pháp tiên tiến hơn như Deep Learning hoặc tăng cường dữ liệu (Data Augmentation) để cải thiện kết quả.
    \item Tăng cường việc chọn lọc và xử lý dữ liệu để giảm nhiễu, nâng cao chất lượng dự đoán.
    \item Nghiên cứu thêm các nguồn dữ liệu thực tế khác để đánh giá khả năng tổng quát hóa của mô hình.
\end{itemize}
\\ Mô hình phân loại dựa trên bộ dữ liệu Breast Cancer Wisconsin đã cho thấy tiềm năng ứng dụng cao trong việc hỗ trợ chẩn đoán ung thư vú, giảm thiểu sai sót và hỗ trợ các bác sĩ trong việc ra quyết định. Tuy nhiên, cần tiếp tục cải tiến và kiểm định trên các dữ liệu mới để đảm bảo tính chính xác và khả năng áp dụng rộng rãi.

